\documentclass[12pt]{article}
\usepackage[utf8]{inputenc}
\usepackage[T1]{fontenc}
\usepackage[margin=1in]{geometry}
\usepackage{setspace}
\usepackage{xcolor}
\usepackage[colorlinks=true, urlcolor=blue!60!black, linkcolor=blue!60!black]{hyperref}
\setlength{\parindent}{0em}
\setlength{\parskip}{0.6em}
\renewcommand{\baselinestretch}{1.0}
\begin{document}
\begin{center}
{\LARGE \textit{United Public Workers v. Mitchell}, 330 U.S. 75 (1947)\par}
\vspace{0.5em}
{\large Case Brief --- Two-Page Edition\par}
\end{center}

\section{Memory Jogger}
Pre-enforcement Hatch Act challenges were nonjusticiable, the Act's ban on partisan management was upheld as to a Mint worker, and Ninth and Tenth Amendment objections fail when Congress has a granted power (\href{https://supreme.justia.com/cases/federal/us/330/75/}{Justia}).

\section{Detailed Case Facts}
The Hatch Act barred most executive-branch employees from taking ``any active part in political management or in political campaigns,'' on pain of removal. Twelve employees and their union sued. Eleven had not violated the Act but wished to campaign off duty. George P. Poole, a U.S. Mint roller, had served as a Democratic ward committeeman, poll worker, and party paymaster, and faced a proposed removal order.

\section{Procedural History}
A three-judge district court found the claims justiciable but upheld the Act. On direct appeal, after reargument, the Supreme Court affirmed in 1947.

\section{Judicial Votes}
4--3 on the merits. Reed wrote for Vinson, Burton, and Frankfurter (who also concurred on jurisdiction). Black and Rutledge dissented, and Douglas dissented in part. Murphy and Jackson did not participate.

\section{Holding}
The eleven non-violators presented no Article III controversy; Poole did. Congress may bar executive employees, even off-duty industrial workers, from active partisan management, and if a granted power exists, Ninth and Tenth Amendment objections ``must fail.''

\section{Analysis of Opinions}
\textbf{Majority:} Courts cannot decide ``abstract'' challenges without ``definite rights'' and ``definite prejudicial interferences.'' Federal powers are ``subtracted from the totality of sovereignty originally in the states and the people.'' The Act extended limits on ``political contributions of money'' to ``political contributions of energy,'' and how to treat industrial workers was a ``matter of detail for Congress.'' \textbf{Frankfurter:} the late docketing defeated jurisdiction. \textbf{Black:} the Act made millions ``mere spectators.'' \textbf{Douglas:} the eleven faced real threats, and a Mint roller is as remote from policy ``as a charwoman.''

\section{Examples: Future Applications}
\textbf{Same-Side:} a ban on a customs inspector running as a party nominee for the state legislature is valid. \textbf{Opposite-Side:} a ban on employees' paid nonpartisan speeches fails under \textit{United States v. NTEU}. \textbf{Fence-Sitter:} an employee's blog urging votes against incumbents without endorsing a party.

\section{Critique}
Randy E. Barnett and Kurt T. Lash argue that Reed's reasoning erases the Ninth Amendment, which \textit{Griswold} later revived. Others see \textit{Mitchell} as a relic of the right-privilege distinction that gave core political speech near-rational-basis deference.

\section{Key Quotations}
\begin{quote}
``The powers granted by the Constitution to the Federal Government are subtracted from the totality of sovereignty originally in the states and the people. Therefore, when objection is made that the exercise of a federal power infringes upon rights reserved by the Ninth and Tenth Amendments, the inquiry must be directed toward the granted power under which the action of the Union was taken. If granted power is found, necessarily the objection of invasion of those rights, reserved by the Ninth and Tenth Amendments, must fail.''
\end{quote}
\noindent --- Justice Reed, for the Court, 330 U.S. at 95--96.

\end{document}
